\documentclass[11pt]{article}
\usepackage{preamble}

\newcommand{\PacketSep}{\ensuremath{\cdot}}
\newcommand{\IconChart}{[CHART]}
\newcommand{\IconClipboard}{[CLIPBOARD]}

\begin{document}

\packetheading{Lesson 5-4 \PacketSep\ Period 3 \PacketSep\ Student Packet}{Radical Equations --- Assessment-Critical: Solve for Variable + Modeling}

\sectionbanner{OBJECTIVES}
{\small
\renewcommand{\arraystretch}{1.25}
\noindent\begin{tabular}{|p{1.8in}|p{4.8in}|}
\hline
\textbf{Math Objective} & Rewrite radical/rational-exponent equations to isolate a specified variable; evaluate such rewritten formulas at given values to produce tables and graphs. \\ \hline
\textbf{Language Objective} & Explain in writing which variable is being isolated and why, using a formula's context (volume, BSA, speed). \\ \hline
\textbf{Essential Understanding} & Solving a formula for a new variable is the SAME algebra as solving a radical equation --- isolate the radical/rational-exponent term, raise to the reciprocal power, clean up. Works for both pure-algebra and real-world cases. \\ \hline
\end{tabular}
}

\sectionbanner{TOPIC GOALS / ESSENTIAL QUESTION / MATERIALS}
{\small
\renewcommand{\arraystretch}{1.25}
\noindent\begin{tabular}{|p{1.8in}|p{4.8in}|}
\hline
\textbf{Topic Goals} & Rewrite formulas to isolate a specified variable (the Topic 5 assessment's core skill); evaluate rewritten formulas at given \(V/x\) values. \\ \hline
\textbf{Essential Question} & How do we extract ANY variable from a radical/rational-exponent formula, and why does the same procedure work for both algebra and real-world contexts? \\ \hline
\textbf{Materials} & Student packet, pencil, calculator for table evaluation. Practice \#44 is the direct assessment rehearsal --- treat it like the real test. \\ \hline
\end{tabular}
}

\sectionbanner{LAUNCH --- Rewrite a Formula + Table Completion (teacher models Ex 2 + \#44)}

\begin{callout}{\IconBook\ ISOLATE-A-VARIABLE RULES}{calloutgreen}
\begin{enumerate}[leftmargin=2.1em,itemsep=2pt,topsep=3pt]
  \item Identify which variable to isolate (look at the final answer form).
  \item Strip away outer operations LAST-IN, FIRST-OUT (addition/subtraction last \(\to\) multiplication/division \(\to\) exponent/radical first to remove).
  \item Raise both sides to the RECIPROCAL power to undo a rational exponent (\(m/n \to n/m\)).
\end{enumerate}
\end{callout}

\begin{callout}{\IconBook\ TABLE-EVALUATION RULES}{calloutyellow}
\begin{enumerate}[leftmargin=2.1em,itemsep=2pt,topsep=3pt]
  \item For each input value in the table, substitute into the rewritten formula.
  \item Simplify rational exponents step-by-step (use calculator for irrational outputs).
  \item Round to specified precision.
  \item The table is a point-cloud of the formula's graph --- plot the points to see the shape.
\end{enumerate}
\end{callout}

\sectionbanner{EXPLORE --- Think-Pair-Share (35 min)}
\textit{Work with a partner. Practice \#10 rehearses Assessment Item 1A (isolate \(x\)). Practice \#44 (done in Launch) rehearses Assessment Item 2 (complete the table). Plan \(\sim\)8 min for \#39 (BSA application).}

\textbf{Bank-item labels (in order):}
\begin{enumerate}[leftmargin=1.6em,itemsep=2pt,topsep=2pt]
  \item Try It 2 --- Rewrite a formula (partner practice)
  \item Practice \#10 --- Isolate \(x\) \IconStar\ Assessment Item 1A rehearsal
  \item Practice \#25 --- Solve for \(y\), rationalize
  \item Practice \#26 --- Solve for \(y\), rationalize
  \item Practice \#27 --- Solve for \(y\), rationalize
  \item Practice \#39 --- BSA application (modeling)
\end{enumerate}

\bankitem{Try It 2 --- Rewrite a formula (partner practice)}{%
The speed, \(v\), of a vehicle in relation to its stopping distance, \(d\), is represented by the equation \(v = 3.57\sqrt{d}\). What is the equation for the stopping distance in terms of the vehicle's speed?
}{}

\bankitem{Practice \#10 --- Isolate \(x\) \IconStar\ Assessment Item 1A rehearsal}{%
Rewrite the equation \(y = \sqrt{\frac{x - 48}{6}}\) to isolate \(x\).
}{}

\bankitem{Practice \#25 --- Solve for \(y\), rationalize}{%
Solve for \(y\). \(x = 3\sqrt[3]{15 + y}\)
}{}

\bankitem{Practice \#26 --- Solve for \(y\), rationalize}{%
Solve for \(y\). \(x = \frac{\sqrt{2y}}{26}\)
}{}

\bankitem{Practice \#27 --- Solve for \(y\), rationalize}{%
Solve for \(y\). \(x = \frac{\sqrt{y - 14.2}}{0.05}\)
}{}

\bankitem{Practice \#39 --- BSA application (modeling)}{%
Solve using the formula \(BSA = \sqrt{\frac{H \cdot M}{3{,}600}}\). A sports medicine specialist determines that a hot-weather training strategy is appropriate for this athlete. To the nearest tenth, what can the mass of the athlete be for the training strategy to be appropriate?

\begin{center}
\begin{tikzpicture}[font=\small, line width=0.9pt]
  \draw[fill=frameshade, rounded corners=4pt] (0,0) rectangle (5.6,2.35);
  \draw[fill=white, rounded corners=3pt] (2.15,0.55) rectangle (3.45,1.95);
  \draw[tealheader] (0.45,0.35) -- (5.15,0.35);
  \draw[tealheader, line width=1.1pt] (4.95,0.15) -- (4.95,2.1);
  \draw[tealheader, line width=1.1pt] (5.1,0.15) -- (5.1,2.1);
  \node[font=\bfseries] at (2.8,1.25) {Athlete};
  \node[anchor=west] at (0.4,2.0) {Finish};
  \node[draw, rounded corners, fill=calloutyellow, anchor=west] at (3.75,1.7) {\(BSA < 2.0\)};
  \node[draw, rounded corners, fill=calloutpeach, anchor=west] at (3.75,1.05) {165 cm};
\end{tikzpicture}
\end{center}
}{}

\sectionbanner{SHARE / SUMMARY (5 min)}

\begin{sentenceframebox}
\textbf{Sentence frames}

Pair share (\#10 or \#39): ``To isolate \_\_\_, I first \_\_\_, then raise both sides to the power \_\_\_. The rewritten formula is \_\_\_. Plugging in \_\_\_ = \_\_\_ gives \_\_\_. In context this means \_\_\_.'' 

Whole class: one pair reads their answer. Class signals thumbs up/sideways.
\end{sentenceframebox}

\begin{callout}{\IconChart\ CONCEPT SUMMARY (your reference for Topic 5 assessment)}{calloutell}
ISOLATING A VARIABLE FROM A RADICAL/RATIONAL-EXPONENT FORMULA:\\
\hspace*{1.4em}--- Identify the variable to isolate and the form of the final answer.\\
\hspace*{1.4em}--- Undo outer operations first (add/sub), then multiply/divide, then remove the exponent/radical.\\
\hspace*{1.4em}--- To undo \(x^{(m/n)}\): raise both sides to the power \((n/m)\).\\
\hspace*{1.4em}--- Simplify: distribute, combine like terms, solve for the target variable.

\medskip
COMPLETING A TABLE FROM A REWRITTEN FORMULA:\\
\hspace*{1.4em}--- Use your rewritten formula (e.g., \(x = \sqrt[3]{4V} / 2\)) as the rule for each row.\\
\hspace*{1.4em}--- Substitute each \(V\)-value. Use a calculator for irrational outputs. Round as directed.\\
\hspace*{1.4em}--- Check: does the table's pattern match the shape you expect (growth/decay/root curve)?\\
\hspace*{1.4em}--- Assessment Item 2 expects THIS process --- from the rewritten formula to the table.
\end{callout}

\sectionbanner{EXIT}

\begin{summaryexitbox}{EXIT TICKET --- Pre-Assessment Confidence Check}
To isolate a variable from a rational-exponent formula I \_\_\_.

When I evaluate the formula at a specific \(x\), I \_\_\_.

One biggest thing I learned today: \_\_\_\_\_\_\_\_\_\_\_\_\_\_\_\_\_\_\_\_\_\_
\end{summaryexitbox}

\clearpage

\sectionbanner{OPTIONAL REINFORCEMENT (not collected)}
\textit{Extra stretch for fast finishers. Try without a calculator first.}

\bankitem{Optional \#1 (Escape Velocity --- performance task)}{%
Performance Task Escape velocity is the velocity at which an object must be traveling to leave a star or planet without falling back to its surface or into orbit. Escape velocity depends on the gravitational constant, \(G = 6.67 \times 10^{-11}\), the mass, \(M\) in kilograms, and radius, \(r\), of the star or planet.

\begin{center}
\begin{tikzpicture}[>=stealth, font=\small, line width=0.9pt]
  \draw[fill=frameshade] (0,0) circle (0.9);
  \node at (0,0) {\textbf{Earth}};
  \draw[tealheader, very thick, ->] (0.85,0.18) .. controls (1.8,1.15) and (3.1,1.85) .. (4.35,1.95);
  \node[draw, rounded corners, fill=calloutyellow, align=center] at (2.45,1.15) {Escape velocity:\\\(v = \sqrt{\frac{2GM}{r}}\)};
  \node[draw, rounded corners, fill=calloutpeach, align=center] at (0,-1.45) {Earth's radius\\\(= 6{,}371{,}000\ \text{m}\)};
\end{tikzpicture}
\end{center}

Part A Rewrite the equation to solve for mass.

Part B Earth has an escape velocity of about 11,200 meters per second. What is Earth's mass in kilograms?
}{}

\end{document}
